\documentclass[11pt,a4paper,oneside]{book}
\usepackage{cmap}
\usepackage[T1]{fontenc}
\usepackage{lmodern}
\usepackage[a4paper,left=1.15in,right=1.15in,top=1.2in,bottom=1.2in]{geometry}
\usepackage{amsmath,amssymb,amsthm,mathtools}
\usepackage{graphicx,booktabs,array,tabularx,longtable,enumitem}
\usepackage{xcolor}
\usepackage{microtype}
\usepackage{fancyhdr}
\usepackage{titlesec}
\usepackage{needspace}
\usepackage{caption}
\usepackage{tikz}
\usetikzlibrary{arrows.meta,positioning,calc,shapes.geometric,decorations.pathreplacing,fit,backgrounds}
\usepackage[hidelinks,bookmarksnumbered,pdfusetitle=false,
  pdftitle={Generative Information Geometry in the UNNS Substrate: Ancestry, Admissibility, Topology and Information Recoverability},
  pdfauthor={UNNS Research Collective},
  pdfsubject={Mathematics of generative histories in the UNNS Substrate: arithmetic admissibility, permutohedral history complexes, forbidden-history duality, information recovery},
  pdfkeywords={UNNS Substrate; generative histories; admissibility; permutohedron; Alexander duality; nerve theorem; information recovery},
  pdfcreator={LaTeX},pdfproducer={pdfLaTeX}]{hyperref}
\usepackage[capitalise,nameinlink]{cleveref}
\usepackage{xurl}
\definecolor{accent}{HTML}{1F5FA8}
\definecolor{ink}{HTML}{222222}
\setlength{\parindent}{1.2em}\setlength{\parskip}{0.15ex}
\renewcommand{\arraystretch}{1.12}
\setlist{itemsep=0.15em,topsep=0.3em}
\sloppy\emergencystretch=2em
\setcounter{tocdepth}{1}
\captionsetup{font=small,labelfont=bf,justification=raggedright,singlelinecheck=false,skip=4pt}
\setlength{\textfloatsep}{12pt plus 2pt minus 2pt}\setlength{\floatsep}{10pt plus 2pt minus 2pt}\setlength{\intextsep}{10pt plus 2pt minus 2pt}
\renewcommand{\topfraction}{0.9}\renewcommand{\bottomfraction}{0.8}\renewcommand{\textfraction}{0.08}\renewcommand{\floatpagefraction}{0.75}
% ---- continuous chapter flow: no forced page break, no blank leaves
\makeatletter
\renewcommand\chapter{\par\addpenalty{-300}\@afterindentfalse\secdef\@chapter\@schapter}
\renewcommand\frontmatter{\pagenumbering{roman}}
\renewcommand\mainmatter{\clearpage\pagenumbering{arabic}}
\renewcommand\backmatter{}
\makeatother
\titleformat{\chapter}[hang]{\normalfont\Large\bfseries\color{ink}}{\thechapter}{0.8em}{}[{\vspace{1pt}\color{accent}\titlerule[0.8pt]}]
\titlespacing*{\chapter}{0pt}{26pt plus 4pt}{14pt plus 2pt}
\titleformat{\section}{\normalfont\large\bfseries}{\thesection}{0.6em}{}
\titlespacing*{\section}{0pt}{14pt plus 3pt minus 2pt}{6pt plus 1pt}
\titleformat{\subsection}{\normalfont\normalsize\bfseries\itshape}{\thesubsection}{0.6em}{}
\titlespacing*{\subsection}{0pt}{10pt plus 2pt minus 1pt}{4pt plus 1pt}
\newcommand{\chapterstart}{\Needspace{0.3\textheight}}
% ---- headers
\pagestyle{fancy}\fancyhf{}
\fancyhead[L]{}
\fancyhead[R]{\small\itshape\nouppercase{\leftmark}}
\fancyfoot[C]{\small\thepage}\renewcommand{\headrulewidth}{0.3pt}
\fancypagestyle{plain}{\fancyhf{}\fancyfoot[C]{\small\thepage}\renewcommand{\headrulewidth}{0pt}}
\renewcommand{\chaptermark}[1]{\markboth{\thechapter\ \ #1}{}}
% ---- theorem environments (one shared counter per chapter)
\theoremstyle{plain}
\newtheorem{theorem}{Theorem}[chapter]
\newtheorem{lemma}[theorem]{Lemma}
\newtheorem{proposition}[theorem]{Proposition}
\newtheorem{corollary}[theorem]{Corollary}
\theoremstyle{definition}
\newtheorem{definition}[theorem]{Definition}
\newtheorem{example}[theorem]{Example}
\newtheorem{observation}[theorem]{Observation}
\theoremstyle{remark}
\newtheorem{remark}[theorem]{Remark}
\crefname{theorem}{Theorem}{Theorems}\crefname{lemma}{Lemma}{Lemmas}\crefname{proposition}{Proposition}{Propositions}
\crefname{corollary}{Corollary}{Corollaries}\crefname{definition}{Definition}{Definitions}\crefname{example}{Example}{Examples}
\crefname{observation}{Observation}{Observations}\crefname{remark}{Remark}{Remarks}
\crefname{figure}{Figure}{Figures}\crefname{table}{Table}{Tables}\crefname{chapter}{Chapter}{Chapters}\crefname{section}{Section}{Sections}
\crefname{appendix}{Appendix}{Appendices}
% ---- scope tags used inside theorem titles (see Appendix A for what each means)
\newcommand{\CA}{computer-assisted}
\newcommand{\FE}{finite-exact}
\newcommand{\BD}{bounded}
\newcommand{\FL}{falsified}
% ---- notation
\newcommand{\F}{\mathbb{F}_2}\newcommand{\Z}{\mathbb{Z}}\newcommand{\N}{\mathbb{N}}\newcommand{\Q}{\mathbb{Q}}
\newcommand{\Hh}{\widetilde{H}}\newcommand{\Pm}{\mathcal{P}}\newcommand{\St}{\mathrm{St}}
\newcommand{\Pref}{\mathrm{Pref}}\newcommand{\lcm}{\operatorname{lcm}}\newcommand{\rk}{\operatorname{rank}}
\newcommand{\Adm}{\mathrm{Adm}}\newcommand{\cl}{\mathrm{cl}}\newcommand{\Pz}{P_z}
\newcommand{\Ent}{\mathrm{H}}\newcommand{\Mut}{\mathrm{I}}\newcommand{\bits}{\text{bits}}
\newcommand{\Kc}{K}\newcommand{\Yc}{Y}\newcommand{\Dd}{D}
\newcommand{\first}{\mathrm{first}}\newcommand{\last}{\mathrm{last}}
\newcommand{\vp}{v}
\newcommand{\code}[1]{\texttt{#1}}
\newcommand{\unns}{\textsc{unns}}
